\chapter{A finite independent-row witness}

The passage from Mack's row-conditioned assumptions to the observed-data filtration is realized
by a concrete model. This section supplies one finite probability space satisfying every hypothesis
used by Theorems~\ref{thm:passage} and~\ref{thm:mack2prime}. The model is a repository witness, while
the assumptions it realizes are Mack's equations (1) and (2) and the filtration $B_k$ used in the
proof of Theorem 2 \cite{Mack1993}.

\begin{definition}[Eight-outcome triangle with row-generated filtration]\label{def:independenceWitness}\lean{VerifiedReserving.IndependenceWitness.Dgen, VerifiedReserving.IndependenceWitness.X, VerifiedReserving.IndependenceWitness.shockEvent}\leanok
\uses{def:RandomTriangle, def:rowAssumptions}
There are three independent sign coordinates on eight equally likely outcomes. Claims in accident
year $i$ are constant at development year zero and depend only on sign $\xi_i$ thereafter. The
filtration at development year $k$ is defined to be the join of the claims in every row through
$k$. The positive-sign event in row $i$ is called \texttt{shockEvent i}.
\end{definition}

\begin{theorem}[Independent rows and the generated filtration]\label{thm:independenceWitnessRows}\lean{VerifiedReserving.IndependenceWitness.iIndepSet_shockEvent, VerifiedReserving.IndependenceWitness.rowsIndep, VerifiedReserving.IndependenceWitness.rowsGenerateD}\leanok
\uses{def:independenceWitness}
The three row shocks are mutually independent. Every claim in row $i$ is measurable with respect
to that row's shock, so the full row sigma-algebras are independent. By construction, the
observed-data filtration is exactly the join of the row histories.
\end{theorem}
\begin{proof}\leanok The eight intersections of the three positive-sign events have the product
probabilities $1$, $1/2$, $1/4$ and $1/8$ according to their cardinality. Independence passes to
the smaller sigma-algebras generated by the claims. The filtration identity is definitional.
\end{proof}

\begin{theorem}[The independence passage is non-vacuous]\label{thm:independenceWitnessPassage}\lean{VerifiedReserving.IndependenceWitness.mack1Row, VerifiedReserving.IndependenceWitness.mack1_from_rows, VerifiedReserving.IndependenceWitness.mack2'_from_rows, VerifiedReserving.IndependenceWitness.nontrivial, VerifiedReserving.IndependenceWitness.exists_independence_witness}\leanok
\uses{thm:independenceWitnessRows, thm:passage, thm:mack2prime}
The model satisfies Mack's row-conditioned mean equation (1). The general independence theorems
therefore give the corresponding mean equation conditioned on $\mathcal{D}_k$ and the cross-row
residual condition (M2'). Two outcomes give different claims in the same accident year, so this is
a genuinely random instance of both conclusions.
\end{theorem}
\begin{proof}\leanok At the first development step, conditional expectation over the trivial row
history averages the two signs. Every later claim is the stated development factor times the
previous claim and is already measurable in the row history. Apply
Theorems~\ref{thm:passage} and~\ref{thm:mack2prime}; direct evaluation at outcomes zero and one
proves nontriviality.
\end{proof}
